\documentclass[11pt,titlepage]{article}

\usepackage{kotex}
\usepackage{amsmath,amssymb,amsthm}
\usepackage{booktabs}
\usepackage{geometry}
\usepackage{hyperref}
\usepackage{graphicx}
\usepackage{caption}
\usepackage{cite}
\usepackage{indentfirst}

\geometry{a4paper, margin=2.5cm}
\hypersetup{hidelinks}

\captionsetup*[figure]{name=Fig., labelsep=period, font=small,
  justification=justified, singlelinecheck=false}
\captionsetup*[table]{name=Table, labelsep=period, font=small,
  justification=centering, singlelinecheck=false}

% Keep only the statement types used by the document.
\theoremstyle{definition}
\newtheorem{definition}{Definition}
\newtheorem{lemma}{Lemma}
\newtheorem{theorem}{Theorem}
\renewcommand{\proofname}{Proof}

\title{보고서 제목:\\
필요한 경우 부제목}
\author{25-059 백재원}
\date{\today}

\begin{document}

\hypersetup{pageanchor=false}
\maketitle
\hypersetup{pageanchor=true}
\pagenumbering{arabic}

% Add an abstract only for a research-style document or when requested.
% \begin{abstract}
% 초록을 작성한다.
% \end{abstract}

% Add a separate Contents page only when it helps navigate a long document.
% \newpage
% \renewcommand{\contentsname}{Contents}
% \tableofcontents
% \newpage

\section{Introduction}

본문을 작성한다.

% Use an external .bib file when references are needed.
% \bibliographystyle{IEEEtran}
% \bibliography{references}

\end{document}
